\section{หลักการและเหตุผล (Principle \& Rationale)}

\hspace*{2em}รูปแบบการยืม-คืนของจากมหาวิทยาลัยในปัจจุบัน มักใช้รูปแบบการจัดเก็บข้อมูลด้วยการจดบันทึกเช่น การบันทึกลงกระดาษ หรือ spreadsheet ที่สะดวกสำหรับการยืมในปริมาณน้อย หรือใช้งานในขณะที่ยังมีอุปกรณ์ในการดูแลไม่มาก ซึ่งทำให้การดูแลและติดตามเป็นไปได้ยากขึ้นเมื่ออุปกรณ์หรือผู้ใช้งานมีจำนวนมากขึ้น

\subsection{สภาพปัญหาในปัจจุบัน (Problem)}

\begin{itemize}
  \item บันทึกข้อมูลด้วยกระดาษ หรือ spreadsheet ทำให้ค้นหา ติดตาม และสรุปข้อมูลได้ยาก และมีโอกาสผิดพลาดสูง
  \item ไม่สามารถตรวจสอบสถานะอุปกรณ์แบบเรียลไทม์ได้ว่าอุปกรณ์ชิ้นใดว่าง ถูกยืม หรืออยู่ระหว่างซ่อม
  \item อุปกรณ์สูญหายหรือส่งคืนล่าช้า เนื่องจากไม่มีระบบแจ้งเตือนและติดตามผู้ยืมอย่างเป็นระบบ
  \item ฐานข้อมูลที่กระจัดกระจายทำให้การจัดซื้อ และจัดสรรอุปกรณ์ขาดประสิทธิภาพ
\end{itemize}

\subsection{แนวทางการพัฒนา (Development Guidelines)}

\hspace*{2em}โครงการต้องการที่จะปรับปรุงรูปแบบการยืม-คืนอุปกรณ์แบบเดิมให้มีประสิทธิภาพมากขึ้น โดยการออกแบบกระบวนการยืม-คืนใหม่ให้ประสานเข้ากับระบบดิจิทอล และปรับปรุงประสิทธิภาพของการจัดการอุปกรณ์ของมหาวิทยาลัย

\subsection{ประโยชน์ที่คาดว่าจะได้รับ (Expected Benefits)}

\begin{itemize}
  \item ลดภาระงานเอกสารและงานซ้ำซ้อนของเจ้าหน้าที่ที่ดูแลอุปกรณ์
  \item เพิ่มความสะดวกและรวดเร็วให้นักศึกษาที่ต้องการขอยืม หรือติดตามสถานะ
  \item เพิ่มความโปร่งใสและความสามารถในการตรวจสอบประวัติย้อนกลับ ของทุกรายการยืม–คืน
  \item ลดอัตราการสูญหายและคืนล่าช้าผ่านระบบแจ้งเตือนอัตโนมัติ
  \item ข้อมูลสถิติเชิงลึกเพื่อสนับสนุนการตัดสินใจในการจัดซื้อหรือดูแลอุปกรณ์
\end{itemize}
